% ----------------------
\subsection{Section 21 (6 April 1830)}
% ----------------------
	\label{DC21}
	
	% -----
	\subsubsection{Historical Background}
	% -----
		\label{DC21-HB}
		
		% --
		\paragraph{JSP}
		% --
			\label{DC21-HB-JSP}
		
			``JS dictated this revelation at the home of Peter Whitmer Sr. in Fayette Township, New York, after the meeting at which the Church of Christ was organized. Such a meeting had been anticipated by JS and his followers for over a year, in part because JS's revelations had foretold the establishment of a church. A March 1829 revelation to Martin Harris, for example, prophesied that a church would be established after three witnesses testified that they had seen the gold plates and after the words found on the plates were `sent forth.' Another JS revelation in spring 1829 confirmed that a church would be established, even as the texts translated from the gold plates prompted JS and Oliver Cowdery to consider questions foundational to a church organization and described how Christ established his church in ancient America. By June, a revelation commanded Cowdery to `build up my church,' and he subsequently drafted a text, which he titled `A commandment from God unto Oliver how he should build up his Church \& the manner thereof.' The document read as though it was intended for immediate implementation.

			``JS, however, later reported that he was directed to delay the establishment of the Church of Christ. His history explained that a revelation `pointed out to us the precise day upon which, according to his will and commandment, we should proceed to organize his Church once again, here upon the earth.' Apparently once a separate June 1829 text, it was incorporated as part of `Articles and Covenants,' a foundational document outlining the governing beliefs, principles, and offices of the church. JS's history described another divine communication, received that June at the home of Peter Whitmer Sr. in Fayette, New York, when the `word of the Lord' instructed JS and Cowdery to wait to ordain each other elders `untill, such times, as it should be practicable' to meet with all the believers who would vote `to accept us as spiritual teachers, or not.' JS and Cowdery were instructed that when the time came they should `then call out such men as the Spirit should dictate, and ordain them, and then attend to the laying on of hands for the gift of the Holy Ghost.'

			``Describing the fulfillment of those instructions, JS's history explained that JS and his associates `met together for that purpose, at the house of...M\supers{r} Whitmer [Peter Whitmer Sr.]' on 6 April 1830. Those who attended the meeting, the history further explained, `consented by an unanimous vote. I then laid my hands upon Oliver Cowdery and ordained him an Elder...after which he ordained me also to the office of an Elder of said Church.' They then partook of the sacramental bread and wine and `then laid...hands on each individual member of the Church present that they might receive the gift of the Holy Ghost, and be confirmed members of the Church of Christ.' Then, while the group was `yet together,' JS dictated the revelation featured below.%
				\footnote{%
					% \label{}%
					All of this ``JS's history'' stuff is from the later history, in volume A-1.
					}

			``It began, `Behold there Shall a Record be kept among you,' and affirmed that in that record JS would be known as a `seer \& Translater \& Prop[h]et an Apostle of Jesus Christ an Elder of the Church.' The revelation instructed `Oliver mine Apostle' to ordain JS. David Whitmer, who was at the meeting, wrote that `Joseph received a revelation that he should be the leader; that he should be ordained by Oliver Cowdery as ``Prophet Seer and Revelator'' to the church, and that the church should receive his words as if from God's own mouth.' When Cowdery was later asked, `To what did you ordain Joseph on the 6th of April, 1830?' he answered, `I ordained him to be a Prophet, Seer, \&c., just as the revelation says.'%
				\footnote{%
					% \label{}%
					JSP note: ``[William E. McLellin], `The Successor of Joseph the Seer,' \textit{Ensign of Liberty}, Dec. 1847, 42.''
					}

			``Eyewitness accounts affirm the date and place this document was dictated; nevertheless, early documents indicate some confusion about both. In the heading to the text itself in Revelation Book 1, the earliest extant copy, and in the table of contents to that volume, John Whitmer dated the revelation to 1829. Later, that date was crossed out multiple times by Oliver Cowdery, who subsequently inserted `April 1830,' after which John Whitmer inserted the number `6' to specify the date as 6 April. Whitmer may have originally written `1829' because the revelation was situated in Revelation Book 1 next to a revelation that does date from 1829. In any case, the text of the revelation supports Cowdery's redaction by referring to the establishment of the church in the past tense.

			``Even though this, the earliest extant version of the revelation, gives Fayette, New York, as the place of dictation, the earliest printed version changed the location to Manchester. The mistake seems to have been the result of confusion on the part of typesetter and printer William W. Phelps in Independence, Missouri, as he prepared this and other revelations for publication in the 1833 Book of Commandments. Phelps, who joined the church in June 1831 and therefore had no firsthand knowledge of the date or place of organization, added `6' to five short April 1830 revelations from Manchester and then followed them with this 6 April text—but he ignored the statement `given at Fayette Seneca County' on the manuscript and instead inserted Manchester to match the others for which he had supplied the same date. Although Phelps and others elsewhere repeated that the church was organized in Manchester, the earliest manuscript and later eyewitness accounts, as well as early financial and legal documents, confirm the correct location as Fayette Township.''
			
	% -----
	\subsubsection{Versions}
	% -----
		\label{DC21-V}
		
		See the \href{https://drive.google.com/file/d/1goAvNz8jS4R8slYfMG_db55Ty2z_vmgK/view?usp=sharing}{sections comparison} document.
	
		\begin{compactdesc}
			\item[JSP] \hyperref[EMS-1830-JS-6-APR-1830]{6 April 1830}
			\item[BC] Chapter 22, 45--46
			\item[DC 1835] Section 46, 177--178
			\item[TS] 3:24 (15 October 1842), 945
			\item[MS] 4:8 (December 1843), 115
			\item[DC 1844] Section 46, 265--266
		\end{compactdesc}

	% -----
	\subsubsection{Section 21:1} % *DC21-1 
	% -----
		\label{DC21-1}

		BEHOLD, there shall be a record kept among you; and in it thou shalt be called a seer, a translator, a prophet, an apostle of Jesus Christ, an elder of the church through the will of God the Father, and the grace of your Lord Jesus Christ,
		
		% \begin{framed}
		% \scriptsize
			% \begin{compactdesc}
				% \item[JSP] 
				% % \item[BC] 
				% % \item[]
			% \end{compactdesc}
		% \end{framed}

	% -----
	\subsubsection{Section 21:2} % *DC21-2 
	% -----
		\label{DC21-2}

		Being inspired of the Holy Ghost to lay the foundation thereof, and to build it up unto the most holy faith.

		% \begin{framed}
		% \scriptsize
			% \begin{compactdesc}
				% \item[JSP] 
				% % \item[BC] 
				% % \item[]
			% \end{compactdesc}
		% \end{framed}

	% -----
	\subsubsection{Section 21:3} % *DC21-3 
	% -----
		\label{DC21-3}

		Which church was organized and established in the year of your Lord eighteen hundred and thirty, in the fourth month, and on the sixth day of the month which is called April.

		% \begin{framed}
		% \scriptsize
			% \begin{compactdesc}
				% \item[JSP] 
				% % \item[BC] 
				% % \item[]
			% \end{compactdesc}
		% \end{framed}

	% -----
	\subsubsection{Section 21:4} % *DC21-4 
	% -----
		\label{DC21-4}

		Wherefore, meaning the church, thou shalt give heed unto all his words and commandments which he shall give unto you as he receiveth them, walking in all holiness before me;

		% \begin{framed}
		% \scriptsize
			% \begin{compactdesc}
				% \item[JSP] 
				% % \item[BC] 
				% % \item[]
			% \end{compactdesc}
		% \end{framed}

	% -----
	\subsubsection{Section 21:5} % *DC21-5 
	% -----
		\label{DC21-5}

		For his word ye shall receive, as if from mine own mouth, in all patience and faith.

		% \begin{framed}
		% \scriptsize
			% \begin{compactdesc}
				% \item[JSP] 
				% % \item[BC] 
				% % \item[]
			% \end{compactdesc}
		% \end{framed}

	% -----
	\subsubsection{Section 21:6} % *DC21-6 
	% -----
		\label{DC21-6}

		For by doing these things the gates of hell shall not prevail against you; yea, and the Lord God will disperse the powers of darkness from before you, and cause the heavens to shake for your good,%
			\footnote{%
				% \label{}%
				See \hyperref[DC35-24]{DC 35:24}.
				}
		and his name's glory.

		% \begin{framed}
		% \scriptsize
			% \begin{compactdesc}
				% \item[JSP] 
				% % \item[BC] 
				% % \item[]
			% \end{compactdesc}
		% \end{framed}

	% -----
	\subsubsection{Section 21:7} % *DC21-7 
	% -----
		\label{DC21-7}

		For thus saith the Lord God: Him have I inspired to move the cause of Zion in mighty power for good, and his diligence I know, and his prayers I have heard.

		% \begin{framed}
		% \scriptsize
			% \begin{compactdesc}
				% \item[JSP] 
				% % \item[BC] 
				% % \item[]
			% \end{compactdesc}
		% \end{framed}

	% -----
	\subsubsection{Section 21:8} % *DC21-8 
	% -----
		\label{DC21-8}

		Yea, his weeping for Zion I have seen, and I will cause that he shall mourn for her no longer; for his days of rejoicing are come unto the remission of his sins, and the manifestations of my blessings upon his works.

		% \begin{framed}
		% \scriptsize
			% \begin{compactdesc}
				% \item[JSP] 
				% % \item[BC] 
				% % \item[]
			% \end{compactdesc}
		% \end{framed}

	% -----
	\subsubsection{Section 21:9} % *DC21-9 
	% -----
		\label{DC21-9}

		For, behold, I will bless all those who labor in my vineyard%
			\footnote{%
				% \label{}%
				This is the first time the DC talks about a ``vineyard.'' It ends up saying a lot about laboring in the vineyard in many different revelations given at different periods.
				}
		with a mighty blessing, and they shall believe on his words, which are given him through me by the \hyperref[COMFORTER]{Comforter}, which manifesteth that Jesus was crucified by sinful men%
			\footnote{%
				% \label{}%
				Some of this language appears to echo \hyperref[LUKE-24-7]{Luke 24:7}.
				}
		for the sins of the world, yea, for the remission of sins unto the \hyperref[CONTRITE]{contrite} heart.

		\begin{framed}
		\scriptsize
			\begin{compactdesc}
				\item[JSP] for behold I will bless all those who Labour in my Vinyard with a mighty blessing \& they shall believe on his words which are given him through me by the comforter which manifesteth that Jesus was Crusified by the Sins of the world for the remision of sins unto the contrite heart
				\item[BC] 10 For behold, I will bless all those who labor in my vineyard, with a mighty blessing, and they shall believe on his words, which are given him through me, by the Comforter: 11 Which manifesteth that Jesus was crucified by sinful men for the sins of the world; 12 Yea, for the remission of sins unto the contrite heart.
				% \item[]
			\end{compactdesc}
		\end{framed}

	% -----
	\subsubsection{Section 21:10} % *DC21-10 
	% -----
		\label{DC21-10}

		Wherefore it behooveth me that he should be ordained by you, Oliver Cowdery mine apostle;%
			\footnote{%
				% \label{}%
				See above, in the \hyperref[DC21-HB-JSP]{historical background}, for what JS is ordained to.
				}

		% \begin{framed}
		% \scriptsize
			% \begin{compactdesc}
				% \item[JSP] 
				% % \item[BC] 
				% % \item[]
			% \end{compactdesc}
		% \end{framed}

	% -----
	\subsubsection{Section 21:11} % *DC21-11 
	% -----
		\label{DC21-11}

		This being an ordinance unto you, that you are an elder under his hand, he being the first unto you, that you might be an elder unto this church of Christ, bearing my name—

		% \begin{framed}
		% \scriptsize
			% \begin{compactdesc}
				% \item[JSP] 
				% % \item[BC] 
				% % \item[]
			% \end{compactdesc}
		% \end{framed}

	% -----
	\subsubsection{Section 21:12} % *DC21-12 
	% -----
		\label{DC21-12}

		And the first preacher of this church unto the church, and before the world, yea, before the Gentiles; yea, and thus saith the Lord God, lo, lo! to the Jews also. Amen.

		% \begin{framed}
		% \scriptsize
			% \begin{compactdesc}
				% \item[JSP] 
				% % \item[BC] 
				% % \item[]
			% \end{compactdesc}
		% \end{framed}